% Created 2026-09-28 Mon 10:58
% Intended LaTeX compiler: pdflatex
\documentclass[11pt]{article}
\usepackage[utf8]{inputenc}
\usepackage[T1]{fontenc}
\usepackage{graphicx}
\usepackage{longtable}
\usepackage{wrapfig}
\usepackage{rotating}
\usepackage[normalem]{ulem}
\usepackage{amsmath}
\usepackage{amssymb}
\usepackage{capt-of}
\usepackage{hyperref}
\author{fcalado}
\date{\today}
\title{}
\hypersetup{
 pdfauthor={fcalado},
 pdftitle={},
 pdfkeywords={},
 pdfsubject={},
 pdfcreator={Emacs 29.3 (Org mode 9.6.15)}, 
 pdflang={English}}
\begin{document}

\tableofcontents

Gender identity is a subjective, internal, malleable and often
transitory sense of oneself"
\begin{itemize}
\item Unlike sex, transgender status is not immutable—-something even a
transgender attorney representing the ACLU had to admit under
questioning from Justice Samuel Alito during oral arguments just two
months ago in United States v. Skrmetti. \textbf{\textbf{Gender identity is a
subjective, internal, malleable and often transitory sense of
oneself}} As such, it is an insufficient basis to expand the notion
of equal protection or protection of “immutable characteristics”—one
that would defy decades of Supreme Court jurisprudence to the
contrary. The congressional letter also cites the court’s decision
in 2020, Bostock v. Clayton County, which held that an employer who
fires an individual merely for being gay or transgender violates
Title VII of the Civil Rights Act---a federal law prohibiting
employment discrimination on the basis of, among other things, sex.
\item "Sorry Democrats, but Trump’s “Two Sexes” Executive Order Is
Constitutional." Sarah Parshall Perry. Mar 20, 2025.
\end{itemize}


We live in a very self-focused subjective psychological world, with
the rise of transgenderism, encouragements to achieve
“self-actualization,” prioritize “me time” and “self-care,” and
exhortations like “pursue your truth.”
\begin{itemize}
\item Expressive individualism is a combination of radical autonomy and
the notion that the inner self is the true self, not necessarily
reflective of or inextricably tied to the physical body. We live in
a very self-focused subjective psychological world, with the rise of
transgenderism, encouragements to achieve “self-actualization,”
prioritize “me time” and “self-care,” and exhortations like “pursue
your truth.” Relationships and morals are viewed as restrictive and
burdensome, rather than ennobling and necessary for human
flourishing. Public policy professor Andrew Cherlin demonstrates
in The Marriage Go-Round that America’s churches, by shifting from a
spirituality of dwelling to a spirituality of seeking, played a role
in promoting expressive individualism.
\item "Feminism: Something To Talk About", Brenda Hafera. Jun 14, 2024
\end{itemize}
\end{document}
